\subsection{A positive-density arithmetic supply for composite conductors}
\label{sec:three-times-artin-prime}

\paragraph{Finding A17 (an explicit family of candidate orders).}
Assume the Dedekind-GRH premise M. Let
\[
 \mathcal Q=\{q\text{ prime}:q\equiv5\pmod8,
                         \operatorname{ord}_q(2)=q-1\}.
\]
Then $\mathcal Q$ has natural density $\mathcal A/2$ among the
primes, where
\[
 \mathcal A=\prod_{\ell\ {
m prime}}
                      \left(1-\frac1{\ell(\ell-1)}\right)>0
\]
is Artin's constant. For every $q\in\mathcal Q$,
\[
                  \langle2,-1\rangle=(\mathbf Z/3q)^\times.
                                                               \tag{AF.1}
\]
The analytic input used here is Lenstra's fixed finite Frobenius
condition version of Hooley's theorem. Its sufficient GRH family is
\[
 \zeta_{K_n}(s),\qquad
 K_n=\mathbf Q(\zeta_n,2^{1/n}),\qquad
 n\ge1\text{ squarefree}.                                    \tag{AF.2}
\]
In particular, this is a number-field Kummer Dedekind-zeta family,
which is contained in M. The density is an arithmetic candidate
supply; it is not a theta-incidence assertion.

\begin{proof}
\emph{1. Degrees and the finite congruence field.}
For odd squarefree $m>1$, the prime $2$ is unramified in
$\mathbf Q(\zeta_m)$. At any prime above $2$, the polynomial
$T^m-2$ is Eisenstein, because the valuation of $2$ is one.
It is irreducible there and hence over $\mathbf Q(\zeta_m)$.
Consequently
\[
                         [K_m:\mathbf Q]=m\varphi(m).
                                                               \tag{AF.3}
\]
The Galois group is the affine semidirect product
$C_m\rtimes(\mathbf Z/m)^\times$. For clarity, a generator
$\tau$ of $C_m$ sends $2^{1/m}$ to $\zeta_m2^{1/m}$, while the
cyclotomic automorphism indexed by $b$ may be chosen to fix
$2^{1/m}$ and sends $\zeta_m$ to $\zeta_m^b$.
Its conjugation sends $\tau$ to $\tau^b$. Since $m$ is odd,
$b=2$ is available, and its commutator with $\tau$ equals
$\tau$. The commutator subgroup therefore contains $C_m$,
and equals it because the quotient by $C_m$ is abelian.
Every abelian subextension of $K_m/\mathbf Q$ thus lies in
$\mathbf Q(\zeta_m)$.

Let $F=\mathbf Q(\zeta_8)$. Its intersection with $K_m$ is an
abelian extension of $\mathbf Q$, so lies in
$\mathbf Q(\zeta_m)\cap\mathbf Q(\zeta_8)=\mathbf Q$.
The latter equality follows from the coprime cyclotomic conductors.
Thus
\[
                       [FK_m:\mathbf Q]=4m\varphi(m).
                                                               \tag{AF.4}
\]
Equations (AF.3)--(AF.4) also hold for $m=1$ with $K_1=\mathbf Q$.

\emph{2. The density from finite inclusion-exclusion.}
The condition $q\equiv5\pmod8$ prescribes the automorphism
$\sigma_5:\zeta_8\mapsto\zeta_8^5$ on $F$. This automorphism
acts nontrivially on $\sqrt2=\zeta_8+\zeta_8^{-1}$, so $2$ is a
quadratic nonresidue at these primes. The prime-$2$ obstruction
to primitive order is already excluded.
For an odd prime $\ell$, the obstruction that
$\ell\mid(q-1)/\operatorname{ord}_q(2)$ is exactly complete
splitting in $K_\ell$, outside the finitely many ramified primes.
For a finite set $E$ of odd primes, the compositum of their
$K_\ell$ is $K_{\prod_{\ell\in E}\ell}$.
By (AF.4), the simultaneous condition of splitting in $K_m$ and
having Frobenius $\sigma_5$ on $F$ has density
$1/(4m\varphi(m))$ for odd squarefree $m$.
Finite inclusion-exclusion therefore gives the density
\[
 \sum_{m\mid\prod_{\ell\in E}\ell}
               \frac{\mu(m)}{4m\varphi(m)}.
                                                               \tag{AF.5}
\]
Lenstra's Theorem~3.1, with $K=\mathbf Q$, $W=\langle2\rangle$,
and the fixed Frobenius condition $(F,\{\sigma_5\})$, identifies
the limiting finite Chebotarev density with the natural density
under (AF.2) \cite[\S2 and Theorem~3.1]{Lenstra1977}.
Thus the analytic infinite-limit step is an application of that
theorem, not an inference from finite inclusion-exclusion alone.
The limit of (AF.5), an absolutely convergent Euler product, is
\[
 \frac14\prod_{\ell\ {
m odd}}
                \left(1-\frac1{\ell(\ell-1)}\right)
                         =\frac{\mathcal A}{2}>0.
\]

\emph{3. The composite unit group.}
For $q\in\mathcal Q$, the Chinese remainder theorem gives
\[
 \operatorname{ord}_{3q}(2)
       =\operatorname{lcm}(2,q-1)=q-1=\frac{\varphi(3q)}2.
\]
A power of $2$ equal to $-1$ modulo $q$ has exponent
$(q-1)/2$ modulo $q-1$, which is even because $q\equiv5\pmod8$.
A power equal to $-1$ modulo $3$ has odd exponent. Therefore
$-1\notin\langle2\rangle$ modulo $3q$. Adjoining it doubles
the subgroup of order $\varphi(3q)/2$, proving (AF.1).
\end{proof}

\paragraph{The exact remaining geometric gate.}
For a fixed genus-two curve over a finite field of
characteristic two, a theta point of exact order $3q$ for even one
$q\in\mathcal Q$ would satisfy A15 on its original locus, and
A18 on every genus-two base. Finding A16 exhibits this for
one particular curve and $q=5$. A17 supplies infinitely many
candidate primes for each fixed curve, but does not establish
such a point for every curve. Poonen's projection theorem with
$S_0\supseteq\{2,3\}$ supplies order $3m$ after prescribing an
exact order-three component; A10 forces $m>1$. Its proof neither
forces $m$ to be a single prime nor places that prime in
$\mathcal Q$. Those additional incidence requirements remain
unproved; the positive density in A17 does not remove them.
